\section{Conclusion}

We presented the first study of nested representation learning for ECG
classification, and subjected its claims to the controls that compression
papers usually omit: matched same-backbone baselines, equivalence testing with
a pre-registered margin, measured rather than extrapolated computation, and
representation probing that could have refuted the mechanism we proposed.

The compressibility result holds. A single Inception1D encoder trained with a
six-level Matryoshka objective attains \R{IncAUCsmall} macro AUC at $d$=16 and
\R{IncAUClarge} at $d$=512 on PTB-XL, differences that are \R{EquivVerdict}
within $\pm$\R{EquivMargin} AUC. One artefact serves every operating point,
with a 32-fold reduction in embedding storage.

The efficiency claim, correctly stated, is narrower than the framing this work
originally carried. Nesting economises on embedding storage, transmission
volume and the number of models an organisation must train, validate and
certify. It does not economise on inference compute---latency is flat in $d$
by construction---and no amount of device-tier tabulation changes that. We
have therefore replaced the wearable-to-cloud framing with the claim the
evidence supports: adaptive embedding compression on clinical twelve-lead ECG,
with reduced-lead and artefact experiments that bound, but do not settle, the
wearable question.

We consider that a more useful contribution than the stronger claim would have
been. Nested representations are a real and practically valuable property of
ECG encoders; establishing precisely where their benefit lies, and where it
does not, is what allows others to build on them correctly.
